\section{Los dispositivos Honghuang: convertir la ruta en una máquina}

Esta sección examina el primer dispositivo de Nengliang Qidian (能量奇点). El valor de Honghuang (洪荒) 70 no consiste en buscar directamente una Q alta, sino en verificar la viabilidad de ingeniería de «construir un dispositivo tokamak completo con materiales superconductores de alta temperatura». Yang Zhao (杨钊) recurrió a una analogía con los barcos: antes todos los barcos se construían con madera, y ahora alguien propone construirlos con acero; el principio de flotación de un barco de acero no cambia, pero sí cambian la soldadura, la eliminación del óxido, la estructura y los riesgos de la botadura. Con el tokamak superconductor de alta temperatura ocurre lo mismo: la física de la fusión no cambia, pero, una vez que cambia el material principal, el diseño, la fabricación, las interfaces y la integración de sistemas deben rehacerse por completo.

Para pasar de una idea a la operación, un dispositivo atraviesa el diseño físico, el diseño conceptual, el diseño de ingeniería, la fabricación, la aceptación de subsistemas, el ensamble general con pruebas conjuntas y la operación experimental. Honghuang 70 comenzó con cuatro fundadores y, al concluir su construcción, contaba con un equipo de alrededor de cien personas; durante el proceso, muchos problemas no se debían a «no saber las fórmulas», sino a no saber en qué etapa estaba el problema, en qué rango caía realmente cada parámetro y cómo decidir cuando no es posible medirlo todo.

\begin{figure}[H]
\centering
\includegraphics[width=0.96\textwidth]{fusion-startup-roadmap.png}
\caption{Hoja de ruta de una empresa emergente de fusión nuclear: de la verificación física a la planta de demostración y, después, a la central comercial. Diagrama conceptual propio, elaborado a partir de la conversación de 01:00:20--01:20:40.}
\end{figure}

\begin{knowledgebox}{Lectura de la figura: la hoja de ruta no es un gran salto}
El paso de Physics a Device y luego a Pilot, Grid y Commercial expresa una reducción de riesgos por etapas. Honghuang 70 verifica la viabilidad de ingeniería de un dispositivo superconductor de alta temperatura completo; el objetivo de Honghuang 170 es lograr Q mayor que 10 con costos controlados; solo Honghuang 380 se asemeja ya a una central de demostración completa y a un producto.
\end{knowledgebox}

\subsection{Qué es el primer plasma}

Esta subsección pasa de la «hoja de ruta» al momento en que «la máquina cobra vida por primera vez». Honghuang 70 no busca alcanzar directamente Q mayor que 10, sino demostrar que un tokamak totalmente superconductor de alta temperatura puede completar las pruebas conjuntas conforme al diseño de ingeniería. Entender el primer plasma evita interpretarlo erróneamente como «ya puede generar electricidad» y también permite comprender por qué sigue siendo un hito de primera importancia.

El primer plasma es la señal de que un dispositivo tokamak realmente «se enciende»: el gas neutro dentro de la cámara de vacío se convierte en plasma mediante preionización, calentamiento y control magnético, y presenta una configuración magnética cercana a la de diseño. No significa que el valor de Q sea alto ni que la central ya sea utilizable; significa que los subsistemas críticos del dispositivo pueden operar en conjunto y, como mínimo, convertir el gas neutro en un plasma controlado.

\begin{importantbox}{El significado del primer plasma}
El primer plasma no es un indicador de generación eléctrica comercial, sino un indicador de integración de sistemas. Muestra que los imanes, el vacío, el calentamiento, la inyección de gas, el diagnóstico y el control cumplen las condiciones básicas de operación coordinada, y que el dispositivo dejó de ser un conjunto de piezas para convertirse en una máquina operable.
\end{importantbox}

\begin{knowledgebox}{Flujo de ingeniería: cómo un tokamak pasa de los planos a la máquina}
\begin{tabularx}{\textwidth}{p{0.18\textwidth}p{0.34\textwidth}X}
\toprule
Etapa & Tareas principales & Riesgo de falla \\
\midrule
Diseño físico & Definir los parámetros objetivo del plasma & Metas demasiado altas o restricciones poco claras provocan desajustes en los sistemas posteriores. \\
Diseño conceptual & Desglosar los indicadores de los sistemas de imanes, vacío, criogenia, calentamiento, diagnóstico, etc. & Los conflictos de interfaz entre subsistemas se amplifican en etapas posteriores. \\
Diseño de ingeniería & Elaborar planos, seleccionar componentes, simular y validar procesos & Los parámetros de las muestras pequeñas no coinciden con los de la pieza real. \\
Fabricación y aceptación & Fabricar los imanes, la cámara de vacío, las fuentes de alimentación y el sistema criogénico, y aceptarlos uno por uno & Degradación del desempeño de piezas individuales o daños durante el transporte o el ensamble. \\
Ensamble general y pruebas conjuntas & Operar varios sistemas a la vez en condiciones de trabajo & Lo más difícil de localizar es «no saber dónde está el problema». \\
\bottomrule
\end{tabularx}
\end{knowledgebox}

\subsection{El imán Jinri y Honghuang 170}

El «imán Jinri (今日磁体)» mencionado en la entrevista es una verificación importante de los imanes de campo toroidal de Honghuang 170. Honghuang 170 debe alcanzar parámetros de campo magnético muy superiores a los de 70, por ejemplo del orden de 23 T; por ello, el equipo primero usó un imán de parámetros altos y gran apertura para verificar el diseño, los procesos y la ruta de fabricación. Este imán no se instalará directamente en 170, pero demuestra que el equipo es capaz de fabricar el subsistema más crítico, más caro y más difícil de los dispositivos futuros.

\begin{knowledgebox}{Descomposición del sistema: dividir un gran riesgo en riesgos de subsistemas}
\begin{tabularx}{\textwidth}{p{0.20\textwidth}p{0.36\textwidth}X}
\toprule
Nivel & Qué se verifica & Problema representativo \\
\midrule
Honghuang 70 & Si un dispositivo superconductor de alta temperatura completo puede construirse y operar & Riesgo de integración de sistemas tras el cambio de generación de materiales. \\
Imán Jinri & Si puede fabricarse un imán de campo alto y gran apertura & Campo magnético alto, alta densidad de corriente de ingeniería, esfuerzos estructurales y enfriamiento criogénico. \\
Honghuang 170 & Si un dispositivo de bajo costo puede lograr Q mayor que 10 & El escalón crítico para pasar de la viabilidad de ingeniería a la viabilidad comercial. \\
Honghuang 380 & Central de demostración y conversión en producto & Operación de larga duración, extracción de calor para generar electricidad, entrega a clientes y rentabilidad de la central. \\
\bottomrule
\end{tabularx}
\end{knowledgebox}

\begin{figure}[H]
\centering
\includegraphics[width=0.96\textwidth]{fusion-bottlenecks.png}
\caption{Cuellos de botella críticos de la fusión nuclear: ciencia, ingeniería, materiales, financiamiento y regulación deben superar el umbral al mismo tiempo. Diagrama conceptual propio, elaborado a partir de la conversación de 01:03:27--01:29:59.}
\end{figure}

\begin{knowledgebox}{Lectura de la figura: los cuellos de botella no se resuelven con un avance aislado}
El plasma, los materiales, los imanes, la ingeniería térmica, el financiamiento y la regulación pueden, cada uno, frenar la comercialización. Al leer la figura conviene atender a la idea de «superar el umbral al mismo tiempo»: que un subsistema sea el primero del mundo no significa que la central sea utilizable, y contar con financiamiento suficiente tampoco significa que desaparezca el riesgo de ingeniería.
\end{knowledgebox}

\subsection{Resumen de la sección}

El significado de Honghuang 70 es convertir el tokamak superconductor de alta temperatura de una idea en un dispositivo completo; el del imán Jinri es verificar la capacidad de construir los subsistemas críticos de 170; Honghuang 170 y 380 corresponden, respectivamente, a lograr Q mayor que 10 y a la comercialización mediante una central de demostración.
